\documentclass[11pt,a4paper]{article}
\usepackage[utf8]{inputenc}
\usepackage[T1]{fontenc}
\usepackage[catalan]{babel}
\usepackage{lmodern}
\usepackage[sfdefault,lining]{FiraSans}
\usepackage[a4paper,top=2.3cm,bottom=2.2cm,left=1.8cm,right=1.8cm,headheight=1.3cm,headsep=0.4cm,footskip=1cm]{geometry}
\usepackage{amsmath,amssymb}
\usepackage{icomma}
\usepackage{graphicx}
\usepackage[table]{xcolor}
\usepackage{tikz}
\usepackage{fontawesome5}
\usepackage{ccicons}
\usepackage{fancyhdr}
\usepackage{lastpage}
\usepackage{tasks}
\usepackage{needspace}
\usepackage[most]{tcolorbox}
\usepackage{array}
\usepackage[hidelinks,pdftitle={Solucions de la fitxa de la unitat 10: Estadística. 1r d'ESO},pdfsubject={Matemàtiques 1r ESO},pdfkeywords={estadística, població, mostra, variables estadístiques, freqüències, taules de freqüències, diagrama de barres, polígon de freqüències, histograma, diagrama de sectors, pictograma, moda, mitjana, mediana, rang, desviació mitjana}]{hyperref}

\definecolor{navy}{HTML}{1B2748}
\definecolor{accent}{HTML}{E8603A}
\definecolor{teal}{HTML}{1F9E90}
\definecolor{soft}{HTML}{EEF1F8}
\definecolor{muted}{HTML}{5D6685}

\setlength{\parindent}{0pt}
\setlength{\parskip}{3pt}
\settasks{label=\alph*), label-format=\bfseries\color{navy}, label-width=1.3em, item-indent=1.8em, column-sep=0.6em, after-item-skip=5pt, before-skip=3pt, after-skip=4pt}

% ---------- capçalera i peu ----------
\newcommand{\llicencia}{\href{https://creativecommons.org/licenses/by-nc-sa/4.0/deed.ca}{\ccbyncsa}\ \ Llicència Creative Commons Reconeixement-NoComercial-CompartirIgual 4.0 Internacional (CC BY-NC-SA 4.0)}
\pagestyle{fancy}
\fancyhf{}
\fancyhead[L]{\raisebox{-0.35\height}{\includegraphics[height=0.95cm]{logo-gva.png}}}
\fancyhead[C]{\small\color{muted}Solucions de la fitxa · Unitat 10 · 1r d'ESO}
\fancyhead[R]{\raisebox{-0.35\height}{\includegraphics[height=0.95cm]{logo-vermellar.png}}}
\fancyfoot[L]{\scriptsize\color{muted}\llicencia}
\fancyfoot[R]{\small\color{muted}\thepage}
\renewcommand{\headrulewidth}{0.4pt}
\renewcommand{\headrule}{\color{muted!40}\hrule height \headrulewidth}
\fancypagestyle{primera}{\fancyhf{}\renewcommand{\headrulewidth}{0pt}
  \fancyfoot[L]{\scriptsize\color{muted}\llicencia}
  \fancyfoot[R]{\small\color{muted}\thepage}}

% ---------- nivells i exercicis ----------
\newcommand{\estrelles}[1]{{\color{accent}\ifcase#1\or\faStar\faStar[regular]\faStar[regular]\or\faStar\faStar\faStar[regular]\or\faStar\faStar\faStar\fi}}
\newcounter{exer}
\newcommand{\ex}[1]{\par\needspace{4\baselineskip}\medskip\stepcounter{exer}\noindent{\color{navy}\bfseries\theexer.}\ \ {\small\estrelles{#1}}\ \ }

\newcommand{\seccio}[2]{\par\needspace{10\baselineskip}\bigskip
  \begin{tcolorbox}[colback=navy,colframe=navy,coltext=white,arc=2mm,boxrule=0pt,left=3mm,right=3mm,top=1.2mm,bottom=1.2mm]
  {\large\bfseries #1\quad #2}
  \end{tcolorbox}\vspace{-2pt}}

\newcolumntype{C}[1]{>{\centering\arraybackslash}p{#1}}
\definecolor{violet}{HTML}{7C6CF0}
\definecolor{gold}{HTML}{E0A21B}

% quadre de recordatori (només a la fitxa d'exercicis)
\newcommand{\recorda}[1]{\par\smallskip\begin{tcolorbox}[colback=accent!7,colframe=accent!55,arc=2mm,boxrule=0.5pt,left=2.5mm,right=2.5mm,top=1mm,bottom=1mm]\small{\color{accent}\faLightbulb[regular]}\ \textbf{Recorda:}\ #1\end{tcolorbox}}
\newcommand{\res}[1]{\boxed{#1}}\newcommand{\rest}[1]{\fbox{#1}}

\def\fmneed{14}
% =====================================================================
% Capçalera FentMates de les fitxes: un color per a cada curs.
% Els colors són els de la pàgina inicial (LEVELS a assets/js/materials.js);
% el logo ve de _eines/logo/pdf/ (python3 _eines/logo/genera.py).
%
% Ús, al final del preàmbul de la unitat:
%   % =====================================================================
% Capçalera FentMates de les fitxes: un color per a cada curs.
% Els colors són els de la pàgina inicial (LEVELS a assets/js/materials.js);
% el logo ve de _eines/logo/pdf/ (python3 _eines/logo/genera.py).
%
% Ús, al final del preàmbul de la unitat:
%   % =====================================================================
% Capçalera FentMates de les fitxes: un color per a cada curs.
% Els colors són els de la pàgina inicial (LEVELS a assets/js/materials.js);
% el logo ve de _eines/logo/pdf/ (python3 _eines/logo/genera.py).
%
% Ús, al final del preàmbul de la unitat:
%   \input{../comu/fentmates.tex}
%   \fmcurs{1eso}
% i, al començament del document:
%   \capcalera{Matemàtiques}{1r ESO}{Fitxa d'exercicis · Unitat 2}{Divisibilitat}
% Opcional abans de \input: \def\fmneed{14} (espai mínim abans de cada secció).
% =====================================================================
\makeatletter
\def\fm@logodir{../../logo/pdf}
\providecommand{\fmneed}{10}
\def\fm@cursnom{}

% clau del curs, color i color del text damunt del color
\newcommand{\fm@curs}[3]{\@namedef{fm@color@#1}{#2}\@namedef{fm@text@#1}{#3}}
\fm@curs{1eso}{E6B84A}{navy}
\fm@curs{2eso}{3B82F6}{white}
\fm@curs{3eso}{FF7A50}{white}
\fm@curs{4eso-a}{A78BFA}{white}
\fm@curs{4eso-b}{6D28D9}{white}
\fm@curs{1batx-mi}{9B1C1C}{white}
\fm@curs{1batx-ccss}{EF6B6B}{white}
\fm@curs{2batx-mii}{1E6B47}{white}
\fm@curs{2batx-ccss}{2BB3A3}{white}

\newcommand{\fmcurs}[1]{%
  \@ifundefined{fm@color@#1}{\PackageError{fentmates}{Curs desconegut: #1}{}}{}%
  \edef\fm@tmp{\noexpand\definecolor{curs}{HTML}{\@nameuse{fm@color@#1}}%
               \noexpand\colorlet{curstext}{\@nameuse{fm@text@#1}}}\fm@tmp
  \colorlet{cursfosc}{curs!72!black}%
  \def\fm@logo{\fm@logodir/vertical-#1.pdf}%
  \def\fm@marca{\fm@logodir/marca-#1.pdf}}

% ---------- capçalera de la primera pàgina ----------
% #1 matèria (petit), #2 curs, #3 tipus de document i unitat, #4 títol
\newcommand{\capcalera}[4]{%
  \gdef\fm@cursnom{#2}%
  \par\noindent
  \begin{tcolorbox}[colback=curs,colframe=curs,boxrule=0pt,arc=2.2mm,width=3.5cm,height=2.15cm,
      valign=center,left=3mm,right=2mm,top=1mm,bottom=1mm,nobeforeafter,box align=center]
    \color{curstext}{\footnotesize #1}\par\vspace{1pt}{\fontsize{19}{21}\selectfont\bfseries #2\par}%
  \end{tcolorbox}\hspace{5mm}%
  \begin{minipage}[c]{\dimexpr\linewidth-3.5cm-5mm-2.9cm\relax}
    {\color{cursfosc}\small\bfseries\MakeUppercase{#3}}\par\vspace{2pt}
    {\color{navy}\fontsize{20}{23}\selectfont\bfseries #4\par}
  \end{minipage}\hfill
  \begin{minipage}[c]{2.6cm}\raggedleft\includegraphics[width=2.6cm]{\fm@logo}\end{minipage}\par
  \vspace{2pt}}

% ---------- barres de secció del color del curs ----------
\renewcommand{\seccio}[2]{\par\needspace{\fmneed\baselineskip}\bigskip
  \begin{tcolorbox}[colback=curs,colframe=curs,coltext=curstext,arc=2mm,boxrule=0pt,left=3mm,right=3mm,top=1.2mm,bottom=1.2mm]
  {\large\bfseries #1\quad #2}
  \end{tcolorbox}\vspace{-2pt}}

% ---------- estrelles de dificultat del color del curs (a Batxillerat no se'n posen) ----------
\@ifundefined{estrelles}{}{%
  \renewcommand{\estrelles}[1]{{\color{curs}\ifcase#1\or\faStar\faStar[regular]\faStar[regular]\or\faStar\faStar\faStar[regular]\or\faStar\faStar\faStar\fi}}}

% ---------- capçalera i peu de les pàgines ----------
% primera pàgina: logos institucionals al peu
\fancypagestyle{primera}{\fancyhf{}\renewcommand{\headrulewidth}{0pt}
  \fancyfoot[L]{\raisebox{-0.32\height}{\includegraphics[height=0.78cm]{logo-gva.png}}\hspace{5mm}%
    \raisebox{-0.32\height}{\includegraphics[height=0.78cm]{logo-vermellar.png}}}
  \fancyfoot[R]{\scriptsize\color{muted}\href{https://creativecommons.org/licenses/by-nc-sa/4.0/deed.ca}{\ccbyncsa}\ \ CC BY-NC-SA 4.0 · fentmates.org\qquad\small\thepage}}
% la resta: el símbol a l'esquerra i el curs a la dreta, amb el títol al centre com abans
\fancyhead[L]{\raisebox{-0.3\height}{\includegraphics[height=0.62cm]{\fm@marca}}}
\fancyhead[R]{\small\bfseries\color{cursfosc}\fm@cursnom}
\renewcommand{\headrule}{\color{curs!45}\hrule height \headrulewidth}
\makeatother

%   \fmcurs{1eso}
% i, al començament del document:
%   \capcalera{Matemàtiques}{1r ESO}{Fitxa d'exercicis · Unitat 2}{Divisibilitat}
% Opcional abans de \input: \def\fmneed{14} (espai mínim abans de cada secció).
% =====================================================================
\makeatletter
\def\fm@logodir{../../logo/pdf}
\providecommand{\fmneed}{10}
\def\fm@cursnom{}

% clau del curs, color i color del text damunt del color
\newcommand{\fm@curs}[3]{\@namedef{fm@color@#1}{#2}\@namedef{fm@text@#1}{#3}}
\fm@curs{1eso}{E6B84A}{navy}
\fm@curs{2eso}{3B82F6}{white}
\fm@curs{3eso}{FF7A50}{white}
\fm@curs{4eso-a}{A78BFA}{white}
\fm@curs{4eso-b}{6D28D9}{white}
\fm@curs{1batx-mi}{9B1C1C}{white}
\fm@curs{1batx-ccss}{EF6B6B}{white}
\fm@curs{2batx-mii}{1E6B47}{white}
\fm@curs{2batx-ccss}{2BB3A3}{white}

\newcommand{\fmcurs}[1]{%
  \@ifundefined{fm@color@#1}{\PackageError{fentmates}{Curs desconegut: #1}{}}{}%
  \edef\fm@tmp{\noexpand\definecolor{curs}{HTML}{\@nameuse{fm@color@#1}}%
               \noexpand\colorlet{curstext}{\@nameuse{fm@text@#1}}}\fm@tmp
  \colorlet{cursfosc}{curs!72!black}%
  \def\fm@logo{\fm@logodir/vertical-#1.pdf}%
  \def\fm@marca{\fm@logodir/marca-#1.pdf}}

% ---------- capçalera de la primera pàgina ----------
% #1 matèria (petit), #2 curs, #3 tipus de document i unitat, #4 títol
\newcommand{\capcalera}[4]{%
  \gdef\fm@cursnom{#2}%
  \par\noindent
  \begin{tcolorbox}[colback=curs,colframe=curs,boxrule=0pt,arc=2.2mm,width=3.5cm,height=2.15cm,
      valign=center,left=3mm,right=2mm,top=1mm,bottom=1mm,nobeforeafter,box align=center]
    \color{curstext}{\footnotesize #1}\par\vspace{1pt}{\fontsize{19}{21}\selectfont\bfseries #2\par}%
  \end{tcolorbox}\hspace{5mm}%
  \begin{minipage}[c]{\dimexpr\linewidth-3.5cm-5mm-2.9cm\relax}
    {\color{cursfosc}\small\bfseries\MakeUppercase{#3}}\par\vspace{2pt}
    {\color{navy}\fontsize{20}{23}\selectfont\bfseries #4\par}
  \end{minipage}\hfill
  \begin{minipage}[c]{2.6cm}\raggedleft\includegraphics[width=2.6cm]{\fm@logo}\end{minipage}\par
  \vspace{2pt}}

% ---------- barres de secció del color del curs ----------
\renewcommand{\seccio}[2]{\par\needspace{\fmneed\baselineskip}\bigskip
  \begin{tcolorbox}[colback=curs,colframe=curs,coltext=curstext,arc=2mm,boxrule=0pt,left=3mm,right=3mm,top=1.2mm,bottom=1.2mm]
  {\large\bfseries #1\quad #2}
  \end{tcolorbox}\vspace{-2pt}}

% ---------- estrelles de dificultat del color del curs (a Batxillerat no se'n posen) ----------
\@ifundefined{estrelles}{}{%
  \renewcommand{\estrelles}[1]{{\color{curs}\ifcase#1\or\faStar\faStar[regular]\faStar[regular]\or\faStar\faStar\faStar[regular]\or\faStar\faStar\faStar\fi}}}

% ---------- capçalera i peu de les pàgines ----------
% primera pàgina: logos institucionals al peu
\fancypagestyle{primera}{\fancyhf{}\renewcommand{\headrulewidth}{0pt}
  \fancyfoot[L]{\raisebox{-0.32\height}{\includegraphics[height=0.78cm]{logo-gva.png}}\hspace{5mm}%
    \raisebox{-0.32\height}{\includegraphics[height=0.78cm]{logo-vermellar.png}}}
  \fancyfoot[R]{\scriptsize\color{muted}\href{https://creativecommons.org/licenses/by-nc-sa/4.0/deed.ca}{\ccbyncsa}\ \ CC BY-NC-SA 4.0 · fentmates.org\qquad\small\thepage}}
% la resta: el símbol a l'esquerra i el curs a la dreta, amb el títol al centre com abans
\fancyhead[L]{\raisebox{-0.3\height}{\includegraphics[height=0.62cm]{\fm@marca}}}
\fancyhead[R]{\small\bfseries\color{cursfosc}\fm@cursnom}
\renewcommand{\headrule}{\color{curs!45}\hrule height \headrulewidth}
\makeatother

%   \fmcurs{1eso}
% i, al començament del document:
%   \capcalera{Matemàtiques}{1r ESO}{Fitxa d'exercicis · Unitat 2}{Divisibilitat}
% Opcional abans de \input: \def\fmneed{14} (espai mínim abans de cada secció).
% =====================================================================
\makeatletter
\def\fm@logodir{../../logo/pdf}
\providecommand{\fmneed}{10}
\def\fm@cursnom{}

% clau del curs, color i color del text damunt del color
\newcommand{\fm@curs}[3]{\@namedef{fm@color@#1}{#2}\@namedef{fm@text@#1}{#3}}
\fm@curs{1eso}{E6B84A}{navy}
\fm@curs{2eso}{3B82F6}{white}
\fm@curs{3eso}{FF7A50}{white}
\fm@curs{4eso-a}{A78BFA}{white}
\fm@curs{4eso-b}{6D28D9}{white}
\fm@curs{1batx-mi}{9B1C1C}{white}
\fm@curs{1batx-ccss}{EF6B6B}{white}
\fm@curs{2batx-mii}{1E6B47}{white}
\fm@curs{2batx-ccss}{2BB3A3}{white}

\newcommand{\fmcurs}[1]{%
  \@ifundefined{fm@color@#1}{\PackageError{fentmates}{Curs desconegut: #1}{}}{}%
  \edef\fm@tmp{\noexpand\definecolor{curs}{HTML}{\@nameuse{fm@color@#1}}%
               \noexpand\colorlet{curstext}{\@nameuse{fm@text@#1}}}\fm@tmp
  \colorlet{cursfosc}{curs!72!black}%
  \def\fm@logo{\fm@logodir/vertical-#1.pdf}%
  \def\fm@marca{\fm@logodir/marca-#1.pdf}}

% ---------- capçalera de la primera pàgina ----------
% #1 matèria (petit), #2 curs, #3 tipus de document i unitat, #4 títol
\newcommand{\capcalera}[4]{%
  \gdef\fm@cursnom{#2}%
  \par\noindent
  \begin{tcolorbox}[colback=curs,colframe=curs,boxrule=0pt,arc=2.2mm,width=3.5cm,height=2.15cm,
      valign=center,left=3mm,right=2mm,top=1mm,bottom=1mm,nobeforeafter,box align=center]
    \color{curstext}{\footnotesize #1}\par\vspace{1pt}{\fontsize{19}{21}\selectfont\bfseries #2\par}%
  \end{tcolorbox}\hspace{5mm}%
  \begin{minipage}[c]{\dimexpr\linewidth-3.5cm-5mm-2.9cm\relax}
    {\color{cursfosc}\small\bfseries\MakeUppercase{#3}}\par\vspace{2pt}
    {\color{navy}\fontsize{20}{23}\selectfont\bfseries #4\par}
  \end{minipage}\hfill
  \begin{minipage}[c]{2.6cm}\raggedleft\includegraphics[width=2.6cm]{\fm@logo}\end{minipage}\par
  \vspace{2pt}}

% ---------- barres de secció del color del curs ----------
\renewcommand{\seccio}[2]{\par\needspace{\fmneed\baselineskip}\bigskip
  \begin{tcolorbox}[colback=curs,colframe=curs,coltext=curstext,arc=2mm,boxrule=0pt,left=3mm,right=3mm,top=1.2mm,bottom=1.2mm]
  {\large\bfseries #1\quad #2}
  \end{tcolorbox}\vspace{-2pt}}

% ---------- estrelles de dificultat del color del curs (a Batxillerat no se'n posen) ----------
\@ifundefined{estrelles}{}{%
  \renewcommand{\estrelles}[1]{{\color{curs}\ifcase#1\or\faStar\faStar[regular]\faStar[regular]\or\faStar\faStar\faStar[regular]\or\faStar\faStar\faStar\fi}}}

% ---------- capçalera i peu de les pàgines ----------
% primera pàgina: logos institucionals al peu
\fancypagestyle{primera}{\fancyhf{}\renewcommand{\headrulewidth}{0pt}
  \fancyfoot[L]{\raisebox{-0.32\height}{\includegraphics[height=0.78cm]{logo-gva.png}}\hspace{5mm}%
    \raisebox{-0.32\height}{\includegraphics[height=0.78cm]{logo-vermellar.png}}}
  \fancyfoot[R]{\scriptsize\color{muted}\href{https://creativecommons.org/licenses/by-nc-sa/4.0/deed.ca}{\ccbyncsa}\ \ CC BY-NC-SA 4.0 · fentmates.org\qquad\small\thepage}}
% la resta: el símbol a l'esquerra i el curs a la dreta, amb el títol al centre com abans
\fancyhead[L]{\raisebox{-0.3\height}{\includegraphics[height=0.62cm]{\fm@marca}}}
\fancyhead[R]{\small\bfseries\color{cursfosc}\fm@cursnom}
\renewcommand{\headrule}{\color{curs!45}\hrule height \headrulewidth}
\makeatother

\fmcurs{1eso}

\begin{document}
\thispagestyle{primera}

\capcalera{Matemàtiques}{1r ESO}{Solucions de la fitxa · Unitat 10}{Estadística}
\vspace{8pt}

\begin{tcolorbox}[colback=soft,colframe=soft,arc=2mm,left=3mm,right=3mm,top=1.5mm,bottom=1.5mm]
\small Intenta resoldre cada exercici abans de mirar-ne la solució. Si no t'ix el mateix resultat, repassa els passos o consulta els apunts interactius de la unitat: hi trobaràs un laboratori de dades que fa taules de freqüències, eines per a construir gràfics i per a calcular la moda, la mitjana, la mediana i la desviació mitjana.
\end{tcolorbox}

% =====================================================
\seccio{1}{Població, mostra i variables estadístiques}

\ex{1} \begin{tasks}(3)
  \task \rest{discreta}
  \task \rest{qualitativa}
  \task \rest{contínua}
  \task \rest{discreta}
  \task \rest{qualitativa}
  \task \rest{contínua}
  \task \rest{discreta}
  \task \rest{qualitativa}
  \task \rest{contínua}
\end{tasks}

\ex{1} Població: \rest{els 450 alumnes de l'institut}. Mostra: \rest{els 60 alumnes enquestats}; grandària \rest{60}. Individu: \rest{cada alumne}. Variable: \rest{hores de televisió setmanals}, quantitativa \rest{contínua}.

\ex{1} \begin{tasks}(4)
  \task \rest{Població}
  \task \rest{Mostra}
  \task \rest{Mostra}
  \task \rest{Població}
\end{tasks}

\ex{2} \begin{tasks}(1)
  \task Per a provar un iogurt cal obrir-lo: si els provàrem tots, \rest{no en quedaria cap per a vendre}.
  \task Són \rest{massa persones}: costaria molts diners i molt de temps.
  \task La prova \rest{destrueix} les bombetes, i a més és lenta i cara.
\end{tasks}

\ex{2} \begin{tasks}(1)
  \task \rest{No}: els qui ixen del gimnàs segurament fan més esport que la resta.
  \task \rest{Sí}: hi ha alumnes de tots els cursos triats a l'atzar.
  \task \rest{No}: els amics de Laura tenen gustos semblants i són de la mateixa edat.
\end{tasks}

\ex{2} \begin{tasks}(3)
  \task 12,5\,kg; 30\,kg; 7,2\,kg: \rest{contínua}
  \task italiana, xinesa, marroquina: \rest{qualitativa}
  \task 0, 3, 8: \rest{discreta}
\end{tasks}

\ex{3} Resposta oberta. Per exemple: població, \rest{els 600 alumnes}; mostra, \rest{6 alumnes triats a l'atzar de cadascun dels 10 grups} (així hi ha alumnes de tots els cursos); pregunta: «Quantes hores llegeixes a la setmana?»; variable quantitativa \rest{contínua}.


% =====================================================
\seccio{2}{Freqüències absolutes i relatives}

\ex{1} \par\smallskip\begin{center}{\renewcommand{\arraystretch}{1.25}\begin{tabular}{|c|c|c|c|}\hline $x_i$ & $f_i$ & $h_i$ & \%\\\hline $0$ & $4$ & $0{,}2$ & $20$\\\hline $1$ & $9$ & $0{,}45$ & $45$\\\hline $2$ & $4$ & $0{,}2$ & $20$\\\hline $3$ & $2$ & $0{,}1$ & $10$\\\hline $4$ & $1$ & $0{,}05$ & $5$\\\hline \textbf{Total} & $20$ & $1$ & $100$\\\hline\end{tabular}}\end{center}

\ex{1} \par\smallskip\begin{center}{\renewcommand{\arraystretch}{1.25}\begin{tabular}{|c|c|c|c|}\hline Grup & $f_i$ & $h_i$ & \%\\\hline A & $9$ & $0{,}36$ & $36$\\\hline O & $11$ & $0{,}44$ & $44$\\\hline B & $3$ & $0{,}12$ & $12$\\\hline AB & $2$ & $0{,}08$ & $8$\\\hline \textbf{Total} & $25$ & $1$ & $100$\\\hline\end{tabular}}\end{center}

\ex{1} \begin{tasks}(4)
  \task $\res{30\,\%}$
  \task $\res{5\,\%}$
  \task $\res{75\,\%}$
  \task $\res{35\,\%}$
\end{tasks}\begin{tasks}(4)
  \task $\res{0{,}12}$
  \task $\res{0{,}6}$
  \task $\res{0{,}075}$
  \task $\res{1}$
\end{tasks}

\ex{2} \par\smallskip\begin{center}{\renewcommand{\arraystretch}{1.25}\begin{tabular}{|c|c|c|c|}\hline Esport & $f_i$ & $h_i$ & \%\\\hline Futbol & $14$ & $0{,}35$ & $35$\\\hline Bàsquet & $10$ & $0{,}25$ & $25$\\\hline Tenis & $6$ & $0{,}15$ & $15$\\\hline Natació & $8$ & $0{,}2$ & $20$\\\hline Altres & $2$ & $0{,}05$ & $5$\\\hline \textbf{Total} & $40$ & $1$ & $100$\\\hline\end{tabular}}\end{center}Bàsquet: $0{,}25\cdot40=10$; tenis: $0{,}15\cdot40=6$; altres: $40-(14+10+6+8)=\res{2}$.

\ex{2} $f_6=60-50=\res{10}$. \ Parells: $11+8+10=29$; \ $\frac{29}{60}\approx\res{48{,}3\,\%}$. \ Més gran que 4: $12+10=22$; \ $\frac{22}{60}\approx\res{36{,}7\,\%}$.

\ex{2} $0{,}35\cdot120=\res{42}$ alumnes. \ $h=\frac{18}{120}=\res{0{,}15}$ (el $15\,\%$).

\ex{3} $0{,}2\cdot N=6\Rightarrow N=6:0{,}2=\res{30}$ alumnes. \ No en porten: $30-6=\res{24}$.

\ex{3} $0{,}4\cdot N=16\Rightarrow N=\res{40}$. \ $h_C=1-0{,}25-0{,}4=\res{0{,}35}$. \ $f_A=0{,}25\cdot40=\res{10}$, \ $f_B=\res{16}$, \ $f_C=0{,}35\cdot40=\res{14}$.


% =====================================================
\seccio{3}{Taules de freqüències}

\ex{1} \par\smallskip\begin{center}{\renewcommand{\arraystretch}{1.25}\begin{tabular}{|c|c|c|c|c|}\hline $x_i$ & $f_i$ & $F_i$ & $h_i$ & \%\\\hline $2$ & $1$ & $1$ & $\frac{1}{30}\approx 0{,}03$ & $\approx 3{,}3$\\\hline $3$ & $2$ & $3$ & $\frac{2}{30}\approx 0{,}07$ & $\approx 6{,}7$\\\hline $4$ & $3$ & $6$ & $0{,}1$ & $10$\\\hline $5$ & $6$ & $12$ & $0{,}2$ & $20$\\\hline $6$ & $7$ & $19$ & $\frac{7}{30}\approx 0{,}23$ & $\approx 23{,}3$\\\hline $7$ & $5$ & $24$ & $\frac{5}{30}\approx 0{,}17$ & $\approx 16{,}7$\\\hline $8$ & $3$ & $27$ & $0{,}1$ & $10$\\\hline $9$ & $2$ & $29$ & $\frac{2}{30}\approx 0{,}07$ & $\approx 6{,}7$\\\hline $10$ & $1$ & $30$ & $\frac{1}{30}\approx 0{,}03$ & $\approx 3{,}3$\\\hline \textbf{Total} & $30$ &  & $1$ & $100$\\\hline\end{tabular}}\end{center}

\ex{1} \begin{tasks}(2)
  \task $F=\res{19}$
  \task $30-6=\res{24}$
  \task $30-24=\res{6}$
  \task $\frac{3}{30}=\res{10\,\%}$
\end{tasks}

\ex{2} Restem cada $F_i$ de l'anterior:\par\smallskip\begin{center}{\renewcommand{\arraystretch}{1.3}\begin{tabular}{|c|C{0.9cm}|C{0.9cm}|C{0.9cm}|C{0.9cm}|C{0.9cm}|C{0.9cm}|}\hline $x_i$ & $0$ & $1$ & $2$ & $3$ & $4$ & $5$\\\hline $f_i$ & $2$ & $3$ & $4$ & $6$ & $3$ & $2$\\\hline\end{tabular}}\end{center}Exactament 3: $f=\res{6}$. \ Almenys 2: $20-5=\res{15}$.

\ex{2} \par\smallskip\begin{center}{\renewcommand{\arraystretch}{1.25}\begin{tabular}{|c|c|c|c|c|}\hline Pes (kg) & Marca & $f_i$ & $F_i$ & \%\\\hline $[35,\ 40)$ & $37{,}5$ & $3$ & $3$ & $12{,}5$\\\hline $[40,\ 45)$ & $42{,}5$ & $5$ & $8$ & $\approx 20{,}8$\\\hline $[45,\ 50)$ & $47{,}5$ & $8$ & $16$ & $\approx 33{,}3$\\\hline $[50,\ 55)$ & $52{,}5$ & $5$ & $21$ & $\approx 20{,}8$\\\hline $[55,\ 60)$ & $57{,}5$ & $3$ & $24$ & $12{,}5$\\\hline \textbf{Total} & & $24$ & & $100$\\\hline\end{tabular}}\end{center}

\ex{2} \begin{tasks}(4)
  \task \rest{$[45,\ 50)$}
  \task \rest{$[40,\ 45)$}
  \task \rest{$[55,\ 60)$}
  \task \rest{$[35,\ 40)$}
\end{tasks}

\ex{3} \par\smallskip\begin{center}{\renewcommand{\arraystretch}{1.25}\begin{tabular}{|c|c|c|c|}\hline Temps (s) & Marca & $f_i$ & $F_i$\\\hline $[12,\ 13)$ & $12{,}5$ & $2$ & $2$\\\hline $[13,\ 14)$ & $13{,}5$ & $5$ & $7$\\\hline $[14,\ 15)$ & $14{,}5$ & $6$ & $13$\\\hline $[15,\ 16)$ & $15{,}5$ & $5$ & $18$\\\hline $[16,\ 17)$ & $16{,}5$ & $2$ & $20$\\\hline\end{tabular}}\end{center}Menys de 14\,s: $F=7$, \ $\frac{7}{20}=\res{35\,\%}$. \ 15\,s o més: $5+2=\res{7}$ alumnes.

\ex{3} $N=F_5=\res{25}$. \ $f_4=F_4-F_3=19-12=\res{7}$. \ Més grans que 3: $25-12=13$, \ $\frac{13}{25}=\res{52\,\%}$.


% =====================================================
\seccio{4}{Diagrames de barres, polígons i histogrames}

\ex{1} \par\smallskip\begin{center}\begin{tikzpicture}[baseline=(current bounding box.center),line cap=round,line join=round]
\draw[muted!30,line width=0.3pt] (0,0.000)--(3.800,0.000); \node[left,font=\tiny,text=muted,inner sep=1.5pt] at (0,0.000) {$0$};
\draw[muted!30,line width=0.3pt] (0,0.467)--(3.800,0.467); \node[left,font=\tiny,text=muted,inner sep=1.5pt] at (0,0.467) {$2$};
\draw[muted!30,line width=0.3pt] (0,0.933)--(3.800,0.933); \node[left,font=\tiny,text=muted,inner sep=1.5pt] at (0,0.933) {$4$};
\draw[muted!30,line width=0.3pt] (0,1.400)--(3.800,1.400); \node[left,font=\tiny,text=muted,inner sep=1.5pt] at (0,1.400) {$6$};
\draw[muted!30,line width=0.3pt] (0,1.867)--(3.800,1.867); \node[left,font=\tiny,text=muted,inner sep=1.5pt] at (0,1.867) {$8$};
\draw[muted!30,line width=0.3pt] (0,2.333)--(3.800,2.333); \node[left,font=\tiny,text=muted,inner sep=1.5pt] at (0,2.333) {$10$};
\draw[muted!30,line width=0.3pt] (0,2.800)--(3.800,2.800); \node[left,font=\tiny,text=muted,inner sep=1.5pt] at (0,2.800) {$12$};
\draw[-latex,muted,line width=0.6pt] (0,0)--(4.100,0) node[right,font=\scriptsize,text=navy]{Grup};
\draw[-latex,muted,line width=0.6pt] (0,0)--(0,3.150) node[above,font=\scriptsize,text=navy]{Persones};
\draw[teal,fill=teal!35] (0.280,0) rectangle (0.820,2.100);
\node[below,font=\scriptsize,inner sep=2pt] at (0.550,0) {A};
\draw[accent,fill=accent!35] (1.180,0) rectangle (1.720,2.567);
\node[below,font=\scriptsize,inner sep=2pt] at (1.450,0) {O};
\draw[violet,fill=violet!35] (2.080,0) rectangle (2.620,0.700);
\node[below,font=\scriptsize,inner sep=2pt] at (2.350,0) {B};
\draw[gold,fill=gold!35] (2.980,0) rectangle (3.520,0.467);
\node[below,font=\scriptsize,inner sep=2pt] at (3.250,0) {AB};
\end{tikzpicture}\end{center}

\ex{1} \par\smallskip\begin{center}{\renewcommand{\arraystretch}{1.3}\begin{tabular}{|c|C{0.9cm}|C{0.9cm}|C{0.9cm}|C{0.9cm}|C{0.9cm}|C{0.9cm}|}\hline $x_i$ & $0$ & $1$ & $2$ & $3$ & $4$ & $5$\\\hline $f_i$ & $3$ & $8$ & $6$ & $4$ & $2$ & $1$\\\hline\end{tabular}}\end{center}$N=\res{24}$ alumnes. \ $Mo=\res{1}$ viatge. \ Alguna vegada: $24-3=\res{21}$.

\ex{1} \par\smallskip\begin{center}\begin{tikzpicture}[baseline=(current bounding box.center),line cap=round,line join=round]
\draw[muted!30,line width=0.3pt] (0,0.000)--(4.700,0.000); \node[left,font=\tiny,text=muted,inner sep=1.5pt] at (0,0.000) {$0$};
\draw[muted!30,line width=0.3pt] (0,0.560)--(4.700,0.560); \node[left,font=\tiny,text=muted,inner sep=1.5pt] at (0,0.560) {$2$};
\draw[muted!30,line width=0.3pt] (0,1.120)--(4.700,1.120); \node[left,font=\tiny,text=muted,inner sep=1.5pt] at (0,1.120) {$4$};
\draw[muted!30,line width=0.3pt] (0,1.680)--(4.700,1.680); \node[left,font=\tiny,text=muted,inner sep=1.5pt] at (0,1.680) {$6$};
\draw[muted!30,line width=0.3pt] (0,2.240)--(4.700,2.240); \node[left,font=\tiny,text=muted,inner sep=1.5pt] at (0,2.240) {$8$};
\draw[muted!30,line width=0.3pt] (0,2.800)--(4.700,2.800); \node[left,font=\tiny,text=muted,inner sep=1.5pt] at (0,2.800) {$10$};
\draw[-latex,muted,line width=0.6pt] (0,0)--(5.000,0) node[right,font=\scriptsize,text=navy]{Mascotes};
\draw[-latex,muted,line width=0.6pt] (0,0)--(0,3.150) node[above,font=\scriptsize,text=navy]{Alumnes};
\draw[teal,fill=teal!35] (0.280,0) rectangle (0.820,1.120);
\node[below,font=\scriptsize,inner sep=2pt] at (0.550,0) {0};
\draw[teal,fill=teal!35] (1.180,0) rectangle (1.720,2.520);
\node[below,font=\scriptsize,inner sep=2pt] at (1.450,0) {1};
\draw[teal,fill=teal!35] (2.080,0) rectangle (2.620,1.120);
\node[below,font=\scriptsize,inner sep=2pt] at (2.350,0) {2};
\draw[teal,fill=teal!35] (2.980,0) rectangle (3.520,0.560);
\node[below,font=\scriptsize,inner sep=2pt] at (3.250,0) {3};
\draw[teal,fill=teal!35] (3.880,0) rectangle (4.420,0.280);
\node[below,font=\scriptsize,inner sep=2pt] at (4.150,0) {4};
\draw[violet,line width=1.1pt] (0.550,1.120)--(1.450,2.520)--(2.350,1.120)--(3.250,0.560)--(4.150,0.280);
\fill[violet] (0.550,1.120) circle (1.6pt);
\fill[violet] (1.450,2.520) circle (1.6pt);
\fill[violet] (2.350,1.120) circle (1.6pt);
\fill[violet] (3.250,0.560) circle (1.6pt);
\fill[violet] (4.150,0.280) circle (1.6pt);
\end{tikzpicture}\end{center}

\ex{2} \par\smallskip\begin{center}\begin{tikzpicture}[baseline=(current bounding box.center),line cap=round,line join=round]
\draw[muted!30,line width=0.3pt] (0,0.000)--(5.200,0.000); \node[left,font=\tiny,text=muted,inner sep=1.5pt] at (0,0.000) {$0$};
\draw[muted!30,line width=0.3pt] (0,0.350)--(5.200,0.350); \node[left,font=\tiny,text=muted,inner sep=1.5pt] at (0,0.350) {$1$};
\draw[muted!30,line width=0.3pt] (0,0.700)--(5.200,0.700); \node[left,font=\tiny,text=muted,inner sep=1.5pt] at (0,0.700) {$2$};
\draw[muted!30,line width=0.3pt] (0,1.050)--(5.200,1.050); \node[left,font=\tiny,text=muted,inner sep=1.5pt] at (0,1.050) {$3$};
\draw[muted!30,line width=0.3pt] (0,1.400)--(5.200,1.400); \node[left,font=\tiny,text=muted,inner sep=1.5pt] at (0,1.400) {$4$};
\draw[muted!30,line width=0.3pt] (0,1.750)--(5.200,1.750); \node[left,font=\tiny,text=muted,inner sep=1.5pt] at (0,1.750) {$5$};
\draw[muted!30,line width=0.3pt] (0,2.100)--(5.200,2.100); \node[left,font=\tiny,text=muted,inner sep=1.5pt] at (0,2.100) {$6$};
\draw[muted!30,line width=0.3pt] (0,2.450)--(5.200,2.450); \node[left,font=\tiny,text=muted,inner sep=1.5pt] at (0,2.450) {$7$};
\draw[muted!30,line width=0.3pt] (0,2.800)--(5.200,2.800); \node[left,font=\tiny,text=muted,inner sep=1.5pt] at (0,2.800) {$8$};
\draw[-latex,muted,line width=0.6pt] (0,0)--(5.500,0) node[right,font=\scriptsize,text=navy]{kg};
\draw[-latex,muted,line width=0.6pt] (0,0)--(0,3.150) node[above,font=\scriptsize,text=navy]{Alumnes};
\draw[teal,fill=teal!30] (0.300,0) rectangle (1.220,1.050);
\draw[teal,fill=teal!30] (1.220,0) rectangle (2.140,1.750);
\draw[teal,fill=teal!30] (2.140,0) rectangle (3.060,2.800);
\draw[teal,fill=teal!30] (3.060,0) rectangle (3.980,1.750);
\draw[teal,fill=teal!30] (3.980,0) rectangle (4.900,1.050);
\node[below,font=\tiny,inner sep=2pt] at (0.300,0) {$35$};
\node[below,font=\tiny,inner sep=2pt] at (1.220,0) {$40$};
\node[below,font=\tiny,inner sep=2pt] at (2.140,0) {$45$};
\node[below,font=\tiny,inner sep=2pt] at (3.060,0) {$50$};
\node[below,font=\tiny,inner sep=2pt] at (3.980,0) {$55$};
\node[below,font=\tiny,inner sep=2pt] at (4.900,0) {$60$};
\draw[violet,line width=1.1pt] (0.000,0.000)--(0.760,1.050)--(1.680,1.750)--(2.600,2.800)--(3.520,1.750)--(4.440,1.050)--(5.200,0.000);
\fill[violet] (0.760,1.050) circle (1.5pt);
\fill[violet] (1.680,1.750) circle (1.5pt);
\fill[violet] (2.600,2.800) circle (1.5pt);
\fill[violet] (3.520,1.750) circle (1.5pt);
\fill[violet] (4.440,1.050) circle (1.5pt);
\end{tikzpicture}\end{center}

\ex{2} \begin{tasks}(1)
  \task \rest{diagrama de barres} (discreta)
  \task \rest{histograma} (contínua, en intervals)
  \task \rest{diagrama de barres} o \rest{de sectors} (qualitativa)
  \task \rest{histograma} (contínua, en intervals)
\end{tasks}

\ex{2} \par\smallskip\begin{center}\begin{tikzpicture}[baseline=(current bounding box.center),line cap=round,line join=round]
\draw[muted!30,line width=0.3pt] (0,0.000)--(6.500,0.000); \node[left,font=\tiny,text=muted,inner sep=1.5pt] at (0,0.000) {$0$};
\draw[muted!30,line width=0.3pt] (0,0.560)--(6.500,0.560); \node[left,font=\tiny,text=muted,inner sep=1.5pt] at (0,0.560) {$2$};
\draw[muted!30,line width=0.3pt] (0,1.120)--(6.500,1.120); \node[left,font=\tiny,text=muted,inner sep=1.5pt] at (0,1.120) {$4$};
\draw[muted!30,line width=0.3pt] (0,1.680)--(6.500,1.680); \node[left,font=\tiny,text=muted,inner sep=1.5pt] at (0,1.680) {$6$};
\draw[muted!30,line width=0.3pt] (0,2.240)--(6.500,2.240); \node[left,font=\tiny,text=muted,inner sep=1.5pt] at (0,2.240) {$8$};
\draw[muted!30,line width=0.3pt] (0,2.800)--(6.500,2.800); \node[left,font=\tiny,text=muted,inner sep=1.5pt] at (0,2.800) {$10$};
\draw[-latex,muted,line width=0.6pt] (0,0)--(6.800,0) node[right,font=\scriptsize,text=navy]{Nota};
\draw[-latex,muted,line width=0.6pt] (0,0)--(0,3.150) node[above,font=\scriptsize,text=navy]{Alumnes};
\node[below,font=\scriptsize,inner sep=2pt] at (0.550,0) {4};
\node[below,font=\scriptsize,inner sep=2pt] at (1.450,0) {5};
\node[below,font=\scriptsize,inner sep=2pt] at (2.350,0) {6};
\node[below,font=\scriptsize,inner sep=2pt] at (3.250,0) {7};
\node[below,font=\scriptsize,inner sep=2pt] at (4.150,0) {8};
\node[below,font=\scriptsize,inner sep=2pt] at (5.050,0) {9};
\node[below,font=\scriptsize,inner sep=2pt] at (5.950,0) {10};
\draw[violet,line width=1.1pt] (0.550,0.280)--(1.450,0.560)--(2.350,1.400)--(3.250,1.960)--(4.150,2.240)--(5.050,1.120)--(5.950,0.840);
\fill[violet] (0.550,0.280) circle (1.6pt);
\fill[violet] (1.450,0.560) circle (1.6pt);
\fill[violet] (2.350,1.400) circle (1.6pt);
\fill[violet] (3.250,1.960) circle (1.6pt);
\fill[violet] (4.150,2.240) circle (1.6pt);
\fill[violet] (5.050,1.120) circle (1.6pt);
\fill[violet] (5.950,0.840) circle (1.6pt);
\draw[accent,line width=1.1pt] (0.550,0.000)--(1.450,0.280)--(2.350,0.840)--(3.250,1.680)--(4.150,2.520)--(5.050,1.960)--(5.950,1.120);
\fill[accent] (0.550,0.000) circle (1.6pt);
\fill[accent] (1.450,0.280) circle (1.6pt);
\fill[accent] (2.350,0.840) circle (1.6pt);
\fill[accent] (3.250,1.680) circle (1.6pt);
\fill[accent] (4.150,2.520) circle (1.6pt);
\fill[accent] (5.050,1.960) circle (1.6pt);
\fill[accent] (5.950,1.120) circle (1.6pt);
\end{tikzpicture}\end{center}El polígon de B (taronja) està desplaçat cap a la dreta: la classe \rest{B} ha tret millors notes.

\ex{3} Si l'eix comença en 90, la barra de C (8 unitats d'altura) sembla \rest{quatre vegades} la de B (2 unitats), però en realitat C només té 6 vots més ($98-92$). Cal començar l'eix en \rest{0}:\par\smallskip\begin{center}\begin{tikzpicture}[baseline=(current bounding box.center),line cap=round,line join=round]
\draw[muted!30,line width=0.3pt] (0,0.000)--(2.900,0.000); \node[left,font=\tiny,text=muted,inner sep=1.5pt] at (0,0.000) {$0$};
\draw[muted!30,line width=0.3pt] (0,0.480)--(2.900,0.480); \node[left,font=\tiny,text=muted,inner sep=1.5pt] at (0,0.480) {$20$};
\draw[muted!30,line width=0.3pt] (0,0.960)--(2.900,0.960); \node[left,font=\tiny,text=muted,inner sep=1.5pt] at (0,0.960) {$40$};
\draw[muted!30,line width=0.3pt] (0,1.440)--(2.900,1.440); \node[left,font=\tiny,text=muted,inner sep=1.5pt] at (0,1.440) {$60$};
\draw[muted!30,line width=0.3pt] (0,1.920)--(2.900,1.920); \node[left,font=\tiny,text=muted,inner sep=1.5pt] at (0,1.920) {$80$};
\draw[muted!30,line width=0.3pt] (0,2.400)--(2.900,2.400); \node[left,font=\tiny,text=muted,inner sep=1.5pt] at (0,2.400) {$100$};
\draw[-latex,muted,line width=0.6pt] (0,0)--(3.200,0) node[right,font=\scriptsize,text=navy]{};
\draw[-latex,muted,line width=0.6pt] (0,0)--(0,2.750) node[above,font=\scriptsize,text=navy]{Vots};
\draw[teal,fill=teal!35] (0.280,0) rectangle (0.820,2.304);
\node[below,font=\scriptsize,inner sep=2pt] at (0.550,0) {A};
\draw[accent,fill=accent!35] (1.180,0) rectangle (1.720,2.208);
\node[below,font=\scriptsize,inner sep=2pt] at (1.450,0) {B};
\draw[violet,fill=violet!35] (2.080,0) rectangle (2.620,2.352);
\node[below,font=\scriptsize,inner sep=2pt] at (2.350,0) {C};
\end{tikzpicture}\end{center}


% =====================================================
\seccio{5}{Diagrames de sectors i pictogrames}

\ex{1} Cada alumne val $360^\circ:36=10^\circ$: \ Futbol $\res{150^\circ}$, Bàsquet $\res{90^\circ}$, Tenis $\res{60^\circ}$, Natació $\res{40^\circ}$, Altres $\res{20^\circ}$.\par\smallskip\begin{center}\begin{tikzpicture}[baseline=(current bounding box.center)]
\fill[teal!70,draw=white,line width=0.8pt] (0,0)--(90.000:1.45) arc[start angle=90.000,end angle=-60.000,radius=1.45cm]--cycle;
\node[font=\scriptsize\bfseries,text=white] at (15.000:0.899) {$150^\circ$};
\fill[accent!70,draw=white,line width=0.8pt] (0,0)--(-60.000:1.45) arc[start angle=-60.000,end angle=-150.000,radius=1.45cm]--cycle;
\node[font=\scriptsize\bfseries,text=white] at (-105.000:0.899) {$90^\circ$};
\fill[violet!70,draw=white,line width=0.8pt] (0,0)--(-150.000:1.45) arc[start angle=-150.000,end angle=-210.000,radius=1.45cm]--cycle;
\node[font=\scriptsize\bfseries,text=white] at (-180.000:0.899) {$60^\circ$};
\fill[gold!70,draw=white,line width=0.8pt] (0,0)--(-210.000:1.45) arc[start angle=-210.000,end angle=-250.000,radius=1.45cm]--cycle;
\node[font=\scriptsize\bfseries,text=white] at (-230.000:0.899) {$40^\circ$};
\fill[navy!70,draw=white,line width=0.8pt] (0,0)--(-250.000:1.45) arc[start angle=-250.000,end angle=-270.000,radius=1.45cm]--cycle;
\fill[teal!70] (1.85,0.590) rectangle ++(0.22,0.22); \node[right,font=\scriptsize] at (2.10,0.700) {Futbol};
\fill[accent!70] (1.85,0.240) rectangle ++(0.22,0.22); \node[right,font=\scriptsize] at (2.10,0.350) {Bàsquet};
\fill[violet!70] (1.85,-0.110) rectangle ++(0.22,0.22); \node[right,font=\scriptsize] at (2.10,0.000) {Tenis};
\fill[gold!70] (1.85,-0.460) rectangle ++(0.22,0.22); \node[right,font=\scriptsize] at (2.10,-0.350) {Natació};
\fill[navy!70] (1.85,-0.810) rectangle ++(0.22,0.22); \node[right,font=\scriptsize] at (2.10,-0.700) {Altres};
\end{tikzpicture}\end{center}

\ex{1} Cada persona val $360^\circ:25=14{,}4^\circ$: \ A: $9\cdot14{,}4=\res{129{,}6^\circ}$, O: $11\cdot14{,}4=\res{158{,}4^\circ}$, B: $3\cdot14{,}4=\res{43{,}2^\circ}$, AB: $2\cdot14{,}4=\res{28{,}8^\circ}$.\par\smallskip\begin{center}\begin{tikzpicture}[baseline=(current bounding box.center)]
\fill[teal!70,draw=white,line width=0.8pt] (0,0)--(90.000:1.45) arc[start angle=90.000,end angle=-39.600,radius=1.45cm]--cycle;
\node[font=\scriptsize\bfseries,text=white] at (25.200:0.899) {$129{,}6^\circ$};
\fill[accent!70,draw=white,line width=0.8pt] (0,0)--(-39.600:1.45) arc[start angle=-39.600,end angle=-198.000,radius=1.45cm]--cycle;
\node[font=\scriptsize\bfseries,text=white] at (-118.800:0.899) {$158{,}4^\circ$};
\fill[violet!70,draw=white,line width=0.8pt] (0,0)--(-198.000:1.45) arc[start angle=-198.000,end angle=-241.200,radius=1.45cm]--cycle;
\node[font=\scriptsize\bfseries,text=white] at (-219.600:0.899) {$43{,}2^\circ$};
\fill[gold!70,draw=white,line width=0.8pt] (0,0)--(-241.200:1.45) arc[start angle=-241.200,end angle=-270.000,radius=1.45cm]--cycle;
\fill[teal!70] (1.85,0.415) rectangle ++(0.22,0.22); \node[right,font=\scriptsize] at (2.10,0.525) {A};
\fill[accent!70] (1.85,0.065) rectangle ++(0.22,0.22); \node[right,font=\scriptsize] at (2.10,0.175) {O};
\fill[violet!70] (1.85,-0.285) rectangle ++(0.22,0.22); \node[right,font=\scriptsize] at (2.10,-0.175) {B};
\fill[gold!70] (1.85,-0.635) rectangle ++(0.22,0.22); \node[right,font=\scriptsize] at (2.10,-0.525) {AB};
\end{tikzpicture}\end{center}

\ex{2} Mòbil: $0{,}4\cdot200=\res{80}$, Consola: $0{,}25\cdot200=\res{50}$, Ordinador: $0{,}2\cdot200=\res{40}$, Tauleta: $0{,}15\cdot200=\res{30}$. \ Comprovem: $80+50+40+30=200$.

\ex{2} Platja: $\frac{100}{360}\cdot72=\res{20}$, Muntanya: $\frac{120}{360}\cdot72=\res{24}$, Ciutat: $\frac{80}{360}\cdot72=\res{16}$, Poble: $\frac{60}{360}\cdot72=\res{12}$.

\ex{2} Gener: $\res{12}$, Febrer: $\res{10}$, Març: $\res{16}$, Abril: $\res{22}$. \ Total: $\res{60}$ bicis.

\ex{2} Cal dibuixar $20:5=\res{4}$, \ $35:5=\res{7}$, \ $15:5=\res{3}$ \ i \ $30:5=\res{6}$ icones.\par\smallskip\begin{center}\begin{tikzpicture}[baseline=(current bounding box.center)]
\node[left,font=\small] at (0,-0.00) {1r A};
\node[text=teal,font=\large] at (0.25,-0.00) {\faShoppingBag};
\node[text=teal,font=\large] at (0.85,-0.00) {\faShoppingBag};
\node[text=teal,font=\large] at (1.45,-0.00) {\faShoppingBag};
\node[text=teal,font=\large] at (2.05,-0.00) {\faShoppingBag};
\node[left,font=\small] at (0,-0.55) {1r B};
\node[text=teal,font=\large] at (0.25,-0.55) {\faShoppingBag};
\node[text=teal,font=\large] at (0.85,-0.55) {\faShoppingBag};
\node[text=teal,font=\large] at (1.45,-0.55) {\faShoppingBag};
\node[text=teal,font=\large] at (2.05,-0.55) {\faShoppingBag};
\node[text=teal,font=\large] at (2.65,-0.55) {\faShoppingBag};
\node[text=teal,font=\large] at (3.25,-0.55) {\faShoppingBag};
\node[text=teal,font=\large] at (3.85,-0.55) {\faShoppingBag};
\node[left,font=\small] at (0,-1.10) {1r C};
\node[text=teal,font=\large] at (0.25,-1.10) {\faShoppingBag};
\node[text=teal,font=\large] at (0.85,-1.10) {\faShoppingBag};
\node[text=teal,font=\large] at (1.45,-1.10) {\faShoppingBag};
\node[left,font=\small] at (0,-1.65) {1r D};
\node[text=teal,font=\large] at (0.25,-1.65) {\faShoppingBag};
\node[text=teal,font=\large] at (0.85,-1.65) {\faShoppingBag};
\node[text=teal,font=\large] at (1.45,-1.65) {\faShoppingBag};
\node[text=teal,font=\large] at (2.05,-1.65) {\faShoppingBag};
\node[text=teal,font=\large] at (2.65,-1.65) {\faShoppingBag};
\node[text=teal,font=\large] at (3.25,-1.65) {\faShoppingBag};
\end{tikzpicture}\end{center}

\ex{3} Cada persona val $72^\circ:18=4^\circ$, \ $N=360:4=\res{90}$ persones. \ Primavera: $0{,}3\cdot90=\res{27}$ persones, \ $27\cdot4=\res{108^\circ}$.


% =====================================================
\seccio{6}{Moda, mitjana i mediana}

\ex{1} $Mo=\res{8}$. \ $\bar x=\frac{46}{7}\res{\approx6{,}57}$. \ Ordenades: $3,\ 5,\ 6,\ \underline{7},\ 8,\ 8,\ 9$: \ $Me=\res{7}$.

\ex{1} $Mo=\res{15}$. \ $\bar x=\frac{80}{6}\res{\approx13{,}33}$. \ Ordenades: $9,\ 11,\ \underline{12,\ 15},\ 15,\ 18$: \ $Me=\frac{12+15}{2}=\res{13{,}5}$.

\ex{1} $Mo=$ \rest{estiu}. La variable és qualitativa: \rest{no} es pot calcular ni la mitjana ni la mediana.

\ex{2} \begin{tasks}(1)
  \task $Mo=\res{9}$; \ $\bar x=\frac{30}{5}=\res{6}$; \ $2,\ 4,\ \underline{6},\ 9,\ 9$: $Me=\res{6}$
  \task Dues modes: $\res{3}$ i $\res{8}$; \ $\bar x=\frac{30}{6}=\res{5}$; \ $Me=\frac{3+7}{2}=\res{5}$
  \task \rest{No hi ha moda} (totes apareixen una vegada); \ $\bar x=\frac{100}{4}=\res{25}$; \ $Me=\frac{20+30}{2}=\res{25}$
\end{tasks}

\ex{2} $Mo=\res{2}$. \ $\bar x=\frac{15}{8}=\res{1{,}875}$ gols. \ Ordenades: $0,\ 1,\ 1,\ \underline{2,\ 2},\ 2,\ 3,\ 4$: \ $Me=\res{2}$.

\ex{2} $\bar x=\frac{1\,700}{6}\res{\approx283{,}33}$ mil €. \ Ordenades: $150,\ 155,\ \underline{160,\ 165},\ 170,\ 900$: \ $Me=\res{162{,}5}$ mil €. \ La \rest{mediana}: la casa de 900 mil € és una dada extrema que estira la mitjana cap amunt.

\ex{3} Per exemple, $\res{4,\ 4,\ 5,\ 8,\ 9}$: suma 30, \ $\bar x=6$; el central és el 5; el 4 és el més repetit. (Hi ha altres solucions.)

\ex{3} $\bar x=\frac{26}{4}=\res{6{,}5}$. \ Amb 5 exàmens, la suma ha de ser $7\cdot5=35$: \ $35-26=\res{9}$.


% =====================================================
\seccio{7}{Paràmetres a partir d'una taula de freqüències}

\ex{1} \par\smallskip\begin{center}{\renewcommand{\arraystretch}{1.25}\begin{tabular}{|c|c|c|c|}\hline $x_i$ & $f_i$ & $F_i$ & $x_i\cdot f_i$\\\hline $0$ & $4$ & $4$ & $0$\\\hline $1$ & $10$ & $14$ & $10$\\\hline $2$ & $6$ & $20$ & $12$\\\hline $3$ & $3$ & $23$ & $9$\\\hline $4$ & $2$ & $25$ & $8$\\\hline \textbf{Total} & $25$ &  & $39$\\\hline\end{tabular}}\end{center}$\bar x=\frac{39}{25}=\res{1{,}56}$. \ La dada central és la 13a: $F=14$ en l'1, \ $Me=\res{1}$. \ $Mo=\res{1}$.

\ex{1} \par\smallskip\begin{center}{\renewcommand{\arraystretch}{1.25}\begin{tabular}{|c|c|c|c|}\hline $x_i$ & $f_i$ & $F_i$ & $x_i\cdot f_i$\\\hline $2$ & $1$ & $1$ & $2$\\\hline $3$ & $2$ & $3$ & $6$\\\hline $4$ & $3$ & $6$ & $12$\\\hline $5$ & $6$ & $12$ & $30$\\\hline $6$ & $7$ & $19$ & $42$\\\hline $7$ & $5$ & $24$ & $35$\\\hline $8$ & $3$ & $27$ & $24$\\\hline $9$ & $2$ & $29$ & $18$\\\hline $10$ & $1$ & $30$ & $10$\\\hline \textbf{Total} & $30$ &  & $179$\\\hline\end{tabular}}\end{center}$\bar x=\frac{179}{30}\res{\approx5{,}97}$. \ Les centrals són la 15a i la 16a: totes dues valen 6 ($F$ passa de 12 a 19), \ $Me=\res{6}$. \ $Mo=\res{6}$ ($f=7$).

\ex{2} \par\smallskip\begin{center}{\renewcommand{\arraystretch}{1.25}\begin{tabular}{|c|c|c|c|}\hline $x_i$ & $f_i$ & $F_i$ & $x_i\cdot f_i$\\\hline $35$ & $2$ & $2$ & $70$\\\hline $36$ & $5$ & $7$ & $180$\\\hline $37$ & $7$ & $14$ & $259$\\\hline $38$ & $8$ & $22$ & $304$\\\hline $39$ & $4$ & $26$ & $156$\\\hline $40$ & $2$ & $28$ & $80$\\\hline \textbf{Total} & $28$ &  & $1049$\\\hline\end{tabular}}\end{center}$\bar x=\frac{1\,049}{28}\res{\approx37{,}46}$. \ Les centrals són la 14a ($37$, perquè $F=14$) i la 15a ($38$): \ $Me=\frac{37+38}{2}=\res{37{,}5}$. \ $Mo=\res{38}$.

\ex{2} \par\smallskip\begin{center}{\renewcommand{\arraystretch}{1.25}\begin{tabular}{|c|c|c|c|}\hline $x_i$ & $f_i$ & $F_i$ & $x_i\cdot f_i$\\\hline $0$ & $5$ & $5$ & $0$\\\hline $1$ & $8$ & $13$ & $8$\\\hline $2$ & $9$ & $22$ & $18$\\\hline $3$ & $6$ & $28$ & $18$\\\hline $4$ & $2$ & $30$ & $8$\\\hline \textbf{Total} & $30$ &  & $52$\\\hline\end{tabular}}\end{center}$\bar x=\frac{52}{30}\res{\approx1{,}73}$. \ Les centrals (15a i 16a) valen 2: $Me=\res{2}$. \ $Mo=\res{2}$.

\ex{2} \par\smallskip\begin{center}{\renewcommand{\arraystretch}{1.25}\begin{tabular}{|c|c|c|}\hline Marca & $f_i$ & $x_i\cdot f_i$\\\hline 37{,}5 & $3$ & \\\hline 42{,}5 & $5$ & \\\hline 47{,}5 & $8$ & \\\hline 52{,}5 & $5$ & \\\hline 57{,}5 & $3$ & \\\hline \textbf{Total} & $24$ & $1\,140$\\\hline\end{tabular}}\end{center}$\bar x\approx\frac{1\,140}{24}=\res{47{,}5}$\,kg.

\ex{3} Sense l'1: $0\cdot4+2\cdot6+3\cdot4=24$ i hi ha 14 dades. \ $\frac{24+a}{14+a}=1{,}5\Rightarrow24+a=21+1{,}5a\Rightarrow3=0{,}5a\Rightarrow\res{a=6}$. \ Comprovem: $\frac{30}{20}=1{,}5$.

\ex{3} \par\smallskip\begin{center}{\renewcommand{\arraystretch}{1.25}\begin{tabular}{|c|c|c|c|}\hline $x_i$ & $f_i$ & $F_i$ & $x_i\cdot f_i$\\\hline $1$ & $2$ & $2$ & $2$\\\hline $2$ & $5$ & $7$ & $10$\\\hline $3$ & $8$ & $15$ & $24$\\\hline $4$ & $3$ & $18$ & $12$\\\hline $5$ & $2$ & $20$ & $10$\\\hline \textbf{Total} & $20$ &  & $58$\\\hline\end{tabular}}\end{center}$\bar x=\frac{58}{20}=\res{2{,}9}$. \ Les centrals (10a i 11a) valen 3: $Me=\res{3}$. \ $Mo=\res{3}$.


% =====================================================
\seccio{8}{Rang i desviació mitjana}

\ex{1} \begin{tasks}(3)
  \task $15-3=\res{12}$
  \task $8-(-4)=\res{12}$
  \task $4{,}8-1{,}3=\res{3{,}5}$
\end{tasks}

\ex{1} $\bar x=\frac{35}{5}=7$. \ Distàncies: $4,\ 2,\ 0,\ 2,\ 4$. \ $DM=\frac{12}{5}=\res{2{,}4}$.

\ex{2} $\bar x=\frac{36}{6}=\res{6}$. \ Rang $=12-2=\res{10}$. \ Distàncies: $4,\ 2,\ 2,\ 0,\ 2,\ 6$; \ $DM=\frac{16}{6}\res{\approx2{,}67}$.

\ex{2} Tots dos tenen $\bar x=\frac{33}{6}=5{,}5$. \ Aina: rang $\res{1}$, \ $DM=\frac{6\cdot0{,}5}{6}=\res{0{,}5}$. \ Biel: rang $9-2=\res{7}$, distàncies $3{,}5;\ 3{,}5;\ 2{,}5;\ 2{,}5;\ 1{,}5;\ 1{,}5$, \ $DM=\frac{15}{6}=\res{2{,}5}$. \ \rest{Aina} és més regular.

\ex{2} Totes dues tenen $\bar x=\res{14}$\,°C. \ X: rang $\res{4}$, \ $DM=\frac{6}{5}=\res{1{,}2}$. \ Y: rang $\res{14}$, \ $DM=\frac{8+6+4+4+2}{5}=\frac{24}{5}=\res{4{,}8}$. \ A la ciutat \rest{Y} les temperatures varien molt més.

\ex{3} $\bar x=\frac{2+6+15+12+10}{15}=\frac{45}{15}=\res{3}$. \ $DM=\frac{2\cdot2+1\cdot3+0\cdot5+1\cdot3+2\cdot2}{15}=\frac{14}{15}\res{\approx0{,}93}$.

\ex{3} \begin{tasks}(1)
  \task \rest{Vertader}: la més gran i la més xicoteta coincideixen.
  \task \rest{Fals}: totes les dades es desplacen igual i les distàncies a la mitjana no canvien.
  \task \rest{Fals}: per exemple, $5,\ 5$ i $0,\ 10$ tenen mitjana 5 però $DM=0$ i $DM=5$.
  \task \rest{Vertader}: és una mitjana de distàncies, que són positives o zero.
\end{tasks}


% =====================================================
\seccio{9}{Problemes d'estadística}

\ex{1} Suma necessària: $6{,}5\cdot4=26$. \ Ja té $6{,}5+7+5{,}5=19$. \ Necessita un $26-19=\res{7}$.

\ex{1} Suma: $12\cdot4=48$. \ $48-(10+15+9)=\res{14}$.

\ex{2} $6{,}4\cdot25=160$; \ $7{,}2\cdot15=108$; \ $\bar x=\frac{160+108}{40}=\frac{268}{40}=\res{6{,}7}$. \ No és la mitjana de les mitjanes perquè les classes \rest{no tenen el mateix nombre d'alumnes}.

\ex{2} \par\smallskip\begin{center}{\renewcommand{\arraystretch}{1.25}\begin{tabular}{|c|c|c|c|}\hline $x_i$ & $f_i$ & $F_i$ & $h_i$\\\hline $0$ & $2$ & $2$ & $0{,}1$\\\hline $1$ & $7$ & $9$ & $0{,}35$\\\hline $2$ & $8$ & $17$ & $0{,}4$\\\hline $3$ & $2$ & $19$ & $0{,}1$\\\hline $4$ & $1$ & $20$ & $0{,}05$\\\hline \textbf{Total} & $20$ &  & $1$\\\hline\end{tabular}}\end{center}\par\smallskip\begin{center}\begin{tikzpicture}[baseline=(current bounding box.center),line cap=round,line join=round]
\draw[muted!30,line width=0.3pt] (0,0.000)--(4.700,0.000); \node[left,font=\tiny,text=muted,inner sep=1.5pt] at (0,0.000) {$0$};
\draw[muted!30,line width=0.3pt] (0,0.350)--(4.700,0.350); \node[left,font=\tiny,text=muted,inner sep=1.5pt] at (0,0.350) {$1$};
\draw[muted!30,line width=0.3pt] (0,0.700)--(4.700,0.700); \node[left,font=\tiny,text=muted,inner sep=1.5pt] at (0,0.700) {$2$};
\draw[muted!30,line width=0.3pt] (0,1.050)--(4.700,1.050); \node[left,font=\tiny,text=muted,inner sep=1.5pt] at (0,1.050) {$3$};
\draw[muted!30,line width=0.3pt] (0,1.400)--(4.700,1.400); \node[left,font=\tiny,text=muted,inner sep=1.5pt] at (0,1.400) {$4$};
\draw[muted!30,line width=0.3pt] (0,1.750)--(4.700,1.750); \node[left,font=\tiny,text=muted,inner sep=1.5pt] at (0,1.750) {$5$};
\draw[muted!30,line width=0.3pt] (0,2.100)--(4.700,2.100); \node[left,font=\tiny,text=muted,inner sep=1.5pt] at (0,2.100) {$6$};
\draw[muted!30,line width=0.3pt] (0,2.450)--(4.700,2.450); \node[left,font=\tiny,text=muted,inner sep=1.5pt] at (0,2.450) {$7$};
\draw[muted!30,line width=0.3pt] (0,2.800)--(4.700,2.800); \node[left,font=\tiny,text=muted,inner sep=1.5pt] at (0,2.800) {$8$};
\draw[-latex,muted,line width=0.6pt] (0,0)--(5.000,0) node[right,font=\scriptsize,text=navy]{Fills};
\draw[-latex,muted,line width=0.6pt] (0,0)--(0,3.150) node[above,font=\scriptsize,text=navy]{Famílies};
\draw[teal,fill=teal!35] (0.280,0) rectangle (0.820,0.700);
\node[below,font=\scriptsize,inner sep=2pt] at (0.550,0) {0};
\draw[teal,fill=teal!35] (1.180,0) rectangle (1.720,2.450);
\node[below,font=\scriptsize,inner sep=2pt] at (1.450,0) {1};
\draw[teal,fill=teal!35] (2.080,0) rectangle (2.620,2.800);
\node[below,font=\scriptsize,inner sep=2pt] at (2.350,0) {2};
\draw[teal,fill=teal!35] (2.980,0) rectangle (3.520,0.700);
\node[below,font=\scriptsize,inner sep=2pt] at (3.250,0) {3};
\draw[teal,fill=teal!35] (3.880,0) rectangle (4.420,0.350);
\node[below,font=\scriptsize,inner sep=2pt] at (4.150,0) {4};
\end{tikzpicture}\end{center}$Mo=\res{2}$; \ $\bar x=\frac{33}{20}=\res{1{,}65}$; \ les centrals (10a i 11a) valen 2: $Me=\res{2}$; \ rang $=4-0=\res{4}$. \ Més d'un fill: $8+2+1=11$, \ $\frac{11}{20}=\res{55\,\%}$. \ Sector: $\frac{8}{20}\cdot360^\circ=\res{144^\circ}$.

\ex{2} $22\cdot11=242$; \ $24\cdot12=288$; \ $288-242=\res{46}$ anys.

\ex{3} Suma errònia: $63$; correcta: $63-9+4=58$; \ $\bar x=\res{5{,}8}$. \ Mediana errònia: $\frac{6+7}{2}=6{,}5$. Correcta (3, 4, 4, 5, 5, 6, 7, 7, 8, 9): $\frac{5+6}{2}=\res{5{,}5}$. Sí que canvia.

\ex{3} \begin{tasks}(1)
  \task \rest{Vertader}: $\bar x=6$ en totes dues.
  \task \rest{Vertader}: la mediana és 7.
  \task \rest{Vertader}: té menys desviació mitjana.
  \task \rest{No es pot saber} amb aquestes dades.
\end{tasks}


\end{document}
